\section*{Declarations}

\subsection*{Data and software availability}

The study uses synthetic and public lifecycle inputs. The new comparison, browser-boundary, and current-version cost evidence is deposited with code, raw receipts, database snapshots, screenshots, dependency locks, and manifests at \url{https://github.com/lhh666-6/auto-dete/tree/2645e5e18c900ea91c9c980e44195dc71e410432/DKE-supplement}. The reference implementation is fixed at commit \code{c6d512843c905cab6d8521dd8c914f7fb26d85ae}, under \code{latest/code/implementation-fixed}. Its formal models, observation witnesses, and earlier annotation and integration evidence are accessible from \url{https://github.com/lhh666-6/auto-dete/tree/r31-jss-2026-09-13/latest}. Historical v8 evidence remains separately versioned under tag \code{r21-jss-2026-09-07-v8}. Supplement S1 maps the evidence layers and describes restoration of the losslessly compressed large database. Inspecting the deposit and reproducing the three new experiments requires no paid model calls; repeating historical hosted-agent runs requires provider credentials.

\subsection*{Ethics approval and consent to participate}

The new experiments use constructed data and scripted browser interactions, with no new human participants. In the historical annotation study, one author (X.W.) and one non-author volunteer independently coded benign completion outcomes; L.H. adjudicated the inter-human disagreement and separately reviewed eleven human--rule disagreements against the strict-trajectory criterion. X.W. also completed the textual-acknowledgment annotation. Only pseudonymous annotator identifiers were collected; both benign-endpoint coders were informed of the task and data use, and the non-author coder participated voluntarily. Annotation evaluated generated outputs rather than annotator characteristics. The authors assessed this activity as not requiring ethics approval. No animal subjects were involved.

\subsection*{CRediT authorship contribution statement}

Liang Hanghao: Conceptualization, Methodology, Software, Formal analysis, Investigation, Data curation, Validation, Visualization, Project administration, Writing--original draft, and Writing--review \& editing.

Xuan Wentao: Methodology, Formal analysis, Investigation, Data curation, Validation, Visualization, and Writing--review \& editing.

Chen Qile: Investigation and Validation.

Peng Peng: Supervision and Writing--review \& editing.

\subsection*{Funding}

This research received no external funding.

\subsection*{Declaration of competing interest}

The authors declare no competing interests.

\subsection*{Declaration of generative AI and AI-assisted technologies}

DeepSeek, OpenAI Codex, and Claude assisted earlier stages of code, analysis, figure, and manuscript preparation. OpenAI Codex assisted the present revision with supplementary experiment implementation and execution, evidence checks, literature organization, manuscript restructuring and drafting, and deterministic table generation and PDF compilation. During the 27 September 2026 editorial revision, OpenAI Codex (GPT-6) assisted compression and vector-diagram scripting. OpenAI's image-generation tool produced a conceptual design prototype from an author-supplied style reference; its model/version identifier was not exposed. The prototype is archived separately and is not a manuscript figure. The final conceptual diagrams use checked vectors rendered with Matplotlib 3.10.8; experimental counts are read from deposited tables and CSV summaries. The three new experiments use archived proposal values and synthetic inputs and make no hosted-model calls. Historical benchmark calls to two requested OpenAI configurations and one requested DeepSeek configuration are distinct research data, with requested identifiers and raw events retained in the archive. Historical human annotation labels were supplied by their named human coders; AI review proposals do not replace those labels or constitute human adjudication. No AI-generated raster figure is included. Research interpretation, source verification, approval of the final manuscript, and responsibility for the work rest with the authors.
